\chapter{Εισαγωγή}
\label{ch:intro}

Το παρόν κεφάλαιο παρουσιάζει το πλαίσιο και το κίνητρο της εργασίας, διατυπώνει το πρόβλημα που εξετάζεται, περιγράφει συνοπτικά την προσέγγιση που ακολουθείται, συνοψίζει τη συνεισφορά της εργασίας και κλείνει με τη δομή των υπόλοιπων κεφαλαίων.

\section{Πλαίσιο και κίνητρο}
\label{sec:intro-motivation}

Τα τελευταία χρόνια, η συνεχής αύξηση της παραγωγής και της διαθεσιμότητας
δεδομένων βίντεο έχει ενισχύσει σημαντικά το ενδιαφέρον για την ανάπτυξη μοντέλων
που μπορούν να κατανοούν τόσο το οπτικό περιεχόμενο όσο και τη χρονική εξέλιξη
μιας σκηνής. Τα μοντέλα αυτά βρίσκουν εφαρμογή σε ένα ευρύ φάσμα εργασιών, όπως
η αναγνώριση ανθρώπινων ενεργειών, η ανίχνευση συμβάντων, η παρακολούθηση
αντικειμένων, η ρομποτική αντίληψη και η αυτόνομη πλοήγηση.

Μεγάλοι όγκοι δεδομένων βίντεο είναι διαθέσιμοι από πηγές όπως το διαδίκτυο,
οι κινητές συσκευές και τα αυτόνομα συστήματα. Ωστόσο, τα περισσότερα από αυτά
τα δεδομένα δεν συνοδεύονται από επισημάνσεις. Η χειροκίνητη
επισήμανση βίντεο αποτελεί ιδιαίτερα χρονοβόρα και δαπανηρή διαδικασία, καθώς
πρέπει να ληφθούν υπόψη τόσο το περιεχόμενο των επιμέρους καρέ όσο και η χρονική
εξέλιξή τους. Επιπλέον, ακόμη και όταν υπάρχουν διαθέσιμες επισημάνσεις, ενδέχεται
να έχουν δημιουργηθεί για διαφορετικό σκοπό και να μην είναι κατάλληλες για την
εκάστοτε εργασία.

Η αυτοεπιβλεπόμενη μάθηση (\tl{self-supervised learning}) επιτρέπει στα
μοντέλα να εκπαιδεύονται αξιοποιώντας πληροφορίες που προκύπτουν αυτόματα από
τα ίδια τα δεδομένα, χωρίς να απαιτείται η χειροκίνητη επισήμανσή τους. Μέσω της αυτοεπιβλεπόμενης προεκπαίδευσης, τα μοντέλα
μπορούν να μαθαίνουν γενικεύσιμες και μεταφέρσιμες οπτικοχρονικές αναπαραστάσεις,
οι οποίες αξιοποιούνται στη συνέχεια σε διαφορετικές εργασίες
(\tl{downstream tasks}).

Στο πλαίσιο αυτό έχουν αναπτυχθεί μοντέλα θεμελίωσης για βίντεο
(\tl{video foundation models}), τα οποία προεκπαιδεύονται σε μεγάλης κλίμακας
σύνολα δεδομένων και μπορούν να αποτελέσουν κοινή βάση για την επίλυση πολλών
εργασιών κατανόησης βίντεο. Χαρακτηριστικό παράδειγμα αποτελεί το \tl{V-JEPA}\cite{adrien2023v-jepa},
το οποίο μαθαίνει αναπαραστάσεις προβλέποντας την πληροφορία αποκρυμμένων
τμημάτων ενός βίντεο στον λανθάνοντα χώρο αναπαράστασης.

Παρά τις υψηλές επιδόσεις τους, τα μοντέλα αυτά συνοδεύονται από σημαντικές
απαιτήσεις σε υπολογιστικούς πόρους. Το μεγάλο πλήθος παραμέτρων αυξάνει τόσο
το κόστος της προεκπαίδευσης όσο και τις απαιτήσεις κατά την εκτέλεση
(\tl{inference}), όπως τον χρόνο επεξεργασίας και την κατανάλωση μνήμης
\tl{GPU} (\tl{VRAM}). Ειδικότερα, το \tl{backbone} του \tl{V-JEPA Large}, το
οποίο χρησιμοποιείται ως μοντέλο αναφοράς στην παρούσα εργασία, διαθέτει περίπου
305 εκατομμύρια παραμέτρους. Οι απαιτήσεις αυτές δυσχεραίνουν την αξιοποίηση του μοντέλου σε εφαρμογές
πραγματικού χρόνου και σε περιβάλλοντα περιορισμένων υπολογιστικών πόρων,
όπως κινητές συσκευές, ενσωματωμένα συστήματα (\tl{embedded systems}) και
αυτόνομα ρομποτικά συστήματα.

Προκύπτει, επομένως, το ερώτημα κατά πόσο μικρότερες αρχιτεκτονικές, με
σημαντικά λιγότερες παραμέτρους και χαμηλότερες υπολογιστικές απαιτήσεις,
μπορούν να επιτύχουν απόδοση συγκρίσιμη με εκείνη ενός μεγάλου μοντέλου
αναπαράστασης βίντεο. Το ερώτημα αυτό αποτελεί το βασικό κίνητρο της παρούσας
διπλωματικής εργασίας.

\section{Στόχος και προσέγγιση της εργασίας}
\label{sec:intro-approach}

Στόχος της παρούσας διπλωματικής εργασίας είναι η διερεύνηση της δυνατότητας
μεταφοράς της γνώσης ενός μεγάλου μοντέλου αναπαράστασης βίντεο, το οποίο έχει
προεκπαιδευτεί μέσω αυτοεπιβλεπόμενης μάθησης, σε μικρότερες αρχιτεκτονικές.
Επιδιώκεται με αυτόν τον τρόπο η μείωση του αριθμού παραμέτρων και των
υπολογιστικών απαιτήσεων, διατηρώντας παράλληλα όσο το δυνατόν μεγαλύτερο
μέρος της απόδοσης του αρχικού μοντέλου.

Για τον σκοπό αυτό αξιοποιείται η Απόσταξη Γνώσης
(\tl{Knowledge Distillation}). Ως μοντέλο-δάσκαλος (\tl{teacher})
χρησιμοποιείται το προεκπαιδευμένο \tl{V-JEPA Large}, σε συνδυασμό με έναν
ταξινομητή τύπου \tl{attentive classifier}, εκπαιδευμένο στο σύνολο δεδομένων
\tl{HMDB51}. Ως μοντέλα-μαθητές (\tl{students}) χρησιμοποιούνται οι μικρότερες
αρχιτεκτονικές \tl{Tiny}, \tl{Small} και \tl{Base}.

Συγκεκριμένα, εξετάζονται οι ακόλουθες στρατηγικές Απόσταξης Γνώσης:
\begin{enumerate}
    \item \textbf{\tl{End-to-End Logit Distillation}}: ολόκληρο το
    μοντέλο-μαθητής εκπαιδεύεται με βάση τις τελικές προβλέψεις
    (\tl{logits}) του μοντέλου-δασκάλου, ώστε να προσεγγίζει την κατανομή
    προβλέψεών του μέσω της απόκλισης \tl{Kullback--Leibler}.

    \item \textbf{\tl{End-to-End Hybrid Feature--Logit Distillation}}: η
    εκπαίδευση ολόκληρου του μοντέλου-μαθητή βασίζεται στον συνδυασμό της
    απόσταξης των τελικών προβλέψεων και της αντιστοίχισης των αναπαραστάσεων
    του μοντέλου-δασκάλου και του μοντέλου-μαθητή.

    \item \textbf{\tl{Backbone Feature Distillation and Attentive Classifier Training}}: 
    αρχικά, το \tl{backbone} του
    μοντέλου-μαθητή εκπαιδεύεται μέσω της αντιστοίχισης των αναπαραστάσεών
    του με εκείνες του μοντέλου-δασκάλου. Στη συνέχεια, το
    \tl{backbone} διατηρείται παγωμένο και εκπαιδεύεται εποπτευόμενα ο
    \tl{attentive classifier}.
\end{enumerate}

Η αξιολόγηση πραγματοποιείται κυρίως στην εργασία αναγνώρισης ανθρώπινων ενεργειών
(\tl{human action recognition}) του συνόλου δεδομένων \tl{HMDB51}, ενώ το
\tl{UCF101} χρησιμοποιείται συμπληρωματικά σε συγκεκριμένα πειράματα. Οι
διαφορετικές στρατηγικές και αρχιτεκτονικές συγκρίνονται τόσο ως προς την
ακρίβεια ταξινόμησης όσο και ως προς την υπολογιστική τους αποδοτικότητα,
λαμβάνοντας υπόψη τον αριθμό παραμέτρων, τον ρυθμό επεξεργασίας (\tl{throughput}) και τη
μέγιστη χρήση μνήμης \tl{GPU}.


\section{Συνεισφορά}
\label{sec:intro-contribution}

Τα μοντέλα θεμελίωσης για βίντεο, όπως το \tl{V-JEPA}, έχουν μελετηθεί
κυρίως ως προς τον τρόπο προεκπαίδευσής τους και την επίδοσή τους σε
εργασίες κατανόησης βίντεο. Αντίθετα, η συμπίεσή τους μέσω απόσταξης
γνώσης σε σημαντικά μικρότερες αρχιτεκτονικές έχει διερευνηθεί σε
μικρότερο βαθμό, καθώς μεγάλο μέρος της σχετικής έρευνας στα δεδομένα
βίντεο αξιοποιεί την απόσταξη ως μηχανισμό εκμάθησης αναπαραστάσεων ή
προεκπαίδευσης νέων μοντέλων. Η συνεισφορά της παρούσας εργασίας
συνοψίζεται στα ακόλουθα σημεία:

\begin{itemize}
    \item \textbf{Διερεύνηση της συμπίεσης του \tl{V-JEPA}:} Εξ όσων
    προέκυψε από τη βιβλιογραφική ανασκόπηση, πρόκειται για την πρώτη
    συστηματική μελέτη της απόσταξης γνώσης από το προεκπαιδευμένο
    \tl{V-JEPA Large} προς τις μικρότερες αρχιτεκτονικές \tl{Base},
    \tl{Small} και \tl{Tiny}, με αντικείμενο την αναγνώριση ανθρώπινων
    ενεργειών.

    \item \textbf{Συγκριτική αξιολόγηση στρατηγικών απόσταξης:}
    Στρατηγικές απόσταξης που έχουν μελετηθεί εκτενώς σε μοντέλα εικόνας
    προσαρμόζονται στα δεδομένα βίντεο και συγκρίνονται σε κοινό
    πειραματικό πλαίσιο, ώστε να εξεταστεί ποιος στόχος απόσταξης είναι
    καταλληλότερος για το \tl{V-JEPA}. Η σύγκριση παρέχει ενδείξεις ότι η
    απόσταξη των αναπαραστάσεων του \tl{backbone} υπερέχει σαφώς της
    απόσταξης μόνο των τελικών προβλέψεων.

    \item \textbf{Απόσταξη με περιορισμένο όγκο δεδομένων:}
    Ενώ στη σχετική βιβλιογραφία η απόσταξη μοντέλων βίντεο εφαρμόζεται
    συχνά σε σύνολα εκατοντάδων χιλιάδων βίντεο, η χρήση τέτοιων συνόλων
    δεν είναι πάντα εφικτή όταν οι υπολογιστικοί πόροι είναι περιορισμένοι.
    Τα αποτελέσματα δείχνουν ότι, ακόμη και με σύνολα της τάξης των
    χιλιάδων βίντεο, μπορεί να διατηρηθεί μεγάλο μέρος της επίδοσης του
    μοντέλου-\tl{teacher}, σε αντίθεση με την εκπαίδευση των ίδιων
    αρχιτεκτονικών από την αρχή, η οποία, με αντίστοιχο όγκο δεδομένων,
    οδήγησε σε πολύ χαμηλή ακρίβεια. Παρέχονται επίσης ενδείξεις ότι τα δεδομένα
    απόσταξης δεν είναι απαραίτητο να προέρχονται από το ίδιο σύνολο με τα
    δεδομένα της τελικής αξιολόγησης, εφόσον έχουν συναφές θεματικό
    περιεχόμενο. Καθώς οι ετικέτες δεν χρησιμοποιούνται κατά την απόσταξη
    του \tl{backbone}, η προσέγγιση μπορεί επιπλέον να αξιοποιήσει πρόσθετα
    μη επισημασμένα βίντεο ως δεδομένα απόσταξης.

    \item \textbf{Αποτίμηση του συμβιβασμού ακρίβειας και υπολογιστικού
    κόστους:} Τα αποσταγμένα μοντέλα αξιολογούνται ως προς την ακρίβεια
    και το υπολογιστικό τους κόστος, αναδεικνύοντας τους διαφορετικούς
    συμβιβασμούς που προσφέρουν οι τρεις αρχιτεκτονικές για εφαρμογές με
    διαφορετικούς περιορισμούς υπολογιστικών πόρων.
\end{itemize}

\section{Δομή της εργασίας}
\label{sec:intro-structure}

Αρχικά, στο Κεφάλαιο~\ref{ch:theory} παρουσιάζεται το θεωρητικό
υπόβαθρο που απαιτείται για την κατανόηση των αρχιτεκτονικών και των
τεχνικών που αξιοποιούνται στην εργασία, με έμφαση στην
αυτοεπιβλεπόμενη μάθηση, το μοντέλο \tl{V-JEPA} και την απόσταξη
γνώσης. Το Κεφάλαιο~\ref{ch:sota} παρουσιάζει συνοπτικά την εξέλιξη των μεθόδων
εκμάθησης αναπαραστάσεων εικόνας και βίντεο, καθώς και σχετικές εργασίες
για την αυτοεπιβλεπόμενη μάθηση, τις τεχνικές συμπίεσης και τις
διαφορετικές προσεγγίσεις απόσταξης γνώσης.

Στο Κεφάλαιο~\ref{ch:datasets} περιγράφονται τα σύνολα δεδομένων, τα
χαρακτηριστικά τους και ο ρόλος τους στα διαφορετικά στάδια της
πειραματικής διαδικασίας. Στη συνέχεια, το
Κεφάλαιο~\ref{ch:implementation} περιγράφει τη διαδικασία φόρτωσης και
προετοιμασίας των δεδομένων, τα μοντέλα αναφοράς, τις τρεις στρατηγικές
απόσταξης γνώσης και τις ρυθμίσεις εκπαίδευσης και υλοποίησης που υιοθετήθηκαν.

Στο Κεφάλαιο~\ref{ch:experiments} παρουσιάζονται και αναλύονται τα
αποτελέσματα των πειραμάτων, ενώ τα τελικά μοντέλα-\tl{student}
συγκρίνονται με το μοντέλο-\tl{teacher} ως προς την ακρίβεια και την
υπολογιστική αποδοτικότητα. Τέλος, στο Κεφάλαιο~\ref{ch:conclusions}
συνοψίζονται τα βασικά ευρήματα της εργασίας, επισημαίνονται οι
περιορισμοί της και προτείνονται κατευθύνσεις για μελλοντική έρευνα.

